\subsection{Pipeline ADC}\label{sec:adc_pipeline}

\subsubsection{Overview}
A chain of \(m\) stages, each a sample and hold, a coarse sub-ADC of
\(n_s\) bits, and an MDAC (Section~\ref{sec:dac_mdac}) forming the
amplified residue for the next stage. All stages work concurrently on
successive samples: throughput of one sample per clock, latency \(m\)
clocks.

\subsubsection{Schematic}
\begin{circuitfig}
  \foreach \i/\lab in {0/Stage 1, 2.6/Stage 2, 5.2/$\ldots$, 7.8/Stage $m$}
    \draw (\i,0) rectangle ++(1.8,1.2) node[midway]{\lab};
  \draw[->] (-1,0.6) node[left]{$V_{IN}$} -- (0,0.6);
  \foreach \i in {1.8, 4.4, 7.0} \draw[->] (\i,0.6) -- ++(0.8,0);
  \foreach \i in {0.9, 3.5, 8.7} \draw[->] (\i,0) -- ++(0,-0.8);
  \draw (0,-1.8) rectangle (9.6,-0.8) node[midway]{delay alignment and digital correction};
  \draw[->] (9.6,-1.3) -- ++(0.8,0) node[right]{$D$};
\end{circuitfig}

\subsubsection{Transfer Function}
\textbf{Large Signal, Time Domain, DC Input Analysis (Q5)}, 1.5 bit stage
(\(b \in \{-1, 0, 1\}\), input range \(\pm V_{REF}\)), the MDAC with \(C_1 = C_2\)
(Section~\ref{sec:dac_mdac}):
\begin{align*}
  b &= \begin{cases} -1 & V_{IN} < -\frac{V_{REF}}{4} \\ 0 & -\frac{V_{REF}}{4} \leq V_{IN} \leq \frac{V_{REF}}{4} \\ +1 & V_{IN} > \frac{V_{REF}}{4} \end{cases} \hspace{1cm}
  V_{RES} = 2 V_{IN} - b V_{REF}
\end{align*}
Range of \(V_{RES}\) in each region:
\begin{align*}
  b = -1:&\ V_{IN} \in [-V_{REF}, -\tfrac{V_{REF}}{4}) \Rightarrow V_{RES} = 2V_{IN} + V_{REF} \in [-V_{REF}, \tfrac{V_{REF}}{2}) \\
  b = 0:&\ V_{IN} \in [-\tfrac{V_{REF}}{4}, \tfrac{V_{REF}}{4}] \Rightarrow V_{RES} = 2V_{IN} \in [-\tfrac{V_{REF}}{2}, \tfrac{V_{REF}}{2}] \\
  b = +1:&\ V_{IN} \in (\tfrac{V_{REF}}{4}, V_{REF}] \Rightarrow V_{RES} = 2V_{IN} - V_{REF} \in (-\tfrac{V_{REF}}{2}, V_{REF}]
\end{align*}
so \(V_{RES} \in [-V_{REF}, V_{REF}]\), the next stage's input range. Unrolling
\(m\) stages (\(V_{RES,0} = V_{IN}\), \(V_{RES,k} = 2V_{RES,k-1} - b_k V_{REF}\)):
\begin{align*}
  V_{RES,1} &= 2V_{IN} - b_1 V_{REF} \\
  V_{RES,2} &= 2(2V_{IN} - b_1 V_{REF}) - b_2 V_{REF} = 2^2 V_{IN} - V_{REF}(b_1 2^1 + b_2 2^0) \\
  V_{RES,m} &= 2^m V_{IN} - V_{REF}\sum_{k=1}^{m} b_k 2^{m-k} \hspace{1cm} \text{(Each step doubles the previous sum and appends } b_k 2^0\text{)} \\
  V_{IN} &= V_{REF}\sum_{k=1}^{m} b_k 2^{-k} + \frac{V_{RES,m}}{2^m} \hspace{1cm} [\div 2^m,\ \text{rearranged}]
\end{align*}
so \(D = \sum_k b_k 2^{m-k}\) (an overlapping add, digits \(b_k\) may be \(-1\))
resolves \(V_{IN}\) to \(\frac{V_{REF}}{2^m}\), plus a final flash on \(V_{RES,m}\).

\subsubsection{Analysis}
\begin{enumerate}
  \item \textbf{Redundancy}: move the upper threshold to \(\frac{V_{REF}}{4} + d\). An input
    just below it still gets \(b = 0\), \(V_{RES} = 2\left(\frac{V_{REF}}{4} + d\right) = \frac{V_{REF}}{2} + 2d\), which stays
    \(\leq V_{REF}\) for \(d \leq \frac{V_{REF}}{4}\) (and symmetrically for \(d < 0\) and the
    lower threshold). The residue stays in range and later stages absorb the
    decision error, so cheap comparators with offsets up to
    \(\pm\frac{V_{REF}}{4}\) suffice.
  \item \textbf{Gain Error}: finite op-amp gain scales the residue by
    \(\frac{AF}{1 + AF}\) (Section~\ref{sec:dac_mdac}):
    \begin{align*}
      V_{RES} &= (2 V_{IN} - b V_{REF})\frac{AF}{1 + AF} = (2 V_{IN} - b V_{REF})\left(1 - \frac{1}{1 + AF}\right) \\
      |V_{RES} - V_{RES,ideal}| &= \frac{|2 V_{IN} - b V_{REF}|}{1 + AF} \leq \frac{V_{REF}}{1 + AF} \approx \frac{V_{REF}}{AF} \hspace{1cm} [|V_{RES,ideal}| \leq V_{REF}]
    \end{align*}
    The back end resolves the remaining \(n - 1\) bits over \(2V_{REF}\), an LSB of
    \(\frac{2V_{REF}}{2^{n-1}} = \frac{V_{REF}}{2^{n-2}}\). Keeping the error below half of it:
    \begin{align*}
      \frac{V_{REF}}{AF} &< \frac{1}{2}\cdot\frac{V_{REF}}{2^{n-2}} = \frac{V_{REF}}{2^{n-1}} \hspace{1cm} \Rightarrow \hspace{1cm}
      AF > 2^{n-1} \hspace{1cm} \Rightarrow \hspace{1cm} A > 2^{n} \ \ [F = \tfrac{1}{2}]
    \end{align*}
    (\(F = \frac{C_2}{C_1 + C_2} = \frac{1}{2}\) for \(C_1 = C_2\), \(C_{in}\) neglected).
  \item \textbf{Noise}: stage \(k\) sits behind a gain of \(2^{k-1}\) (stages
    \(1 \ldots k-1\)), so its noise referred to the input is divided by
    \(2^{k-1}\) (power by \(4^{k-1}\)). The first stage dominates and later stages
    are scaled down in capacitance and power.
\end{enumerate}

\subsubsection{Question}
12 bit pipeline of 1.5 bit stages, \(C_1 = C_2\). Find the minimum op-amp DC
gain of the first stage.
\paragraph{Solution:}
\begin{align*}
  A &> 2^{n} = 2^{12} = 4096 \hspace{1cm} (20\log_{10} 4096 = 72.2\,dB)
\end{align*}
